% Deutsche Parallelfassung zu foundations/ellipsoid_geometry.tex

\chapter{Auf der Erde realisierte Krümmung}
\label{chapter:ellipsoid_curvature}

Eine ebene Ingenieurrealisierung behandelt das anwendbare Gebiet als
euklidische Ebene. Für viele Entwurfs-, Bau- und Instandhaltungsaufgaben ist
dies nützlich und hinreichend genau. Ihre Grenze wird sichtbar, wenn die
Ingenieurfolge wesentlich von Projektionsverzerrung oder der Erdkrümmung selbst
abhängt. Eine projizierte CAD-Kurve bleibt dann eine wertvolle
Koordinatenrepräsentation, kann aber weder die intrinsische Konstruktion
besitzen noch die Flächenmetrik bestimmen, aus der ihre Koordinaten abgeleitet
wurden.

Die Frage dieses Kapitels ist daher genau: Was ändert sich, wenn dasselbe
identifizierte Alignment einschließlich der in \cref{chap:state-model}
eingeführten Krümmungsänderung auf der Erde statt in einem ebenen
Ingenieurraum realisiert wird?

\section{Was bleibt und was sich ändert}

Alignment-Identität, intrinsische longitudinale Ordnung, Provenienz der
Änderung und ingenieurtechnische Rolle der Krümmung bleiben invariant. Der
Realization Context ändert sich. Statt von \(\mathcal G_c\) eine ebene
metrische Trajektorie zu verlangen, lassen wir \(\mathcal G_{\mathcal E}\)
den Zustand auf einer Referenzfläche
\[
\mathcal E\subset\mathbb R^3
\]
interpretieren. Das Ergebnis ist eine Flächenkurve
\[
\gamma_{\mathcal E}:[0,L]\rightarrow\mathcal E,
\qquad \lVert\gamma_{\mathcal E}'(s)\rVert=1.
\]
Damit behält \(s\) seine longitudinale Bedeutung, während Position, Richtung
und Krümmung eine Erdoberflächeninterpretation erhalten.

\begin{figure}[htbp]
\centering
\input{figures/foundations/earth_realization_decomposition-DE}
\caption{Die Erdrealisierung bewahrt das intrinsische Alignment und trennt
seine Krümmung innerhalb der Fläche von der durch die Referenzfläche
erzwungenen Krümmung. Die Projektion erzeugt eine weitere Repräsentation der
realisierten Kurve, keine neue konstruktive Definition.}
\label{fig:earth-realization-decomposition}
\end{figure}

\section{Warum eine Tangente nicht mehr genügt}

Die realisierte Tangente
\[
\mathbf t(s)=\gamma_{\mathcal E}'(s)
\]
gibt die Fahrtrichtung an. Auf einer Ebene bleibt eine feste Normale dieser
Ebene üblicherweise implizit. Auf dem Ellipsoid ändert sich die
Flächennormale \(\mathbf n_{\mathcal E}(s)\) mit der Position und muss daher
ausdrücklich in die Realisierung eingehen. Die orientierte Seitenrichtung in
der Tangentialebene ist
\[
\mathbf b(s)=\mathbf n_{\mathcal E}(s)\times\mathbf t(s).
\]
Diese Richtungen werden eingeführt, weil der Krümmungsvektor nun zwei
unterscheidbare Ursachen besitzt.

\section{Die Krümmungszerlegung}

Die Änderung der Tangente im umgebenden Raum zerfällt als
\[
\gamma_{\mathcal E}''(s)
=
\underbrace{\kappa_g(s)\,\mathbf b(s)}_{
\text{Drehung innerhalb der Fläche}}
+
\underbrace{\kappa_n(s)\,\mathbf n_{\mathcal E}(s)}_{
\text{durch die Fläche erzwungene Krümmung}},
\]
wobei für die gewählte Orientierung
\[
\kappa_g(s)=
\left\langle\nabla_s\mathbf t(s),\mathbf b(s)\right\rangle,
\qquad
\kappa_n(s)=
\left\langle\gamma_{\mathcal E}''(s),
\mathbf n_{\mathcal E}(s)\right\rangle.
\]
Folglich gilt
\[
\lVert\gamma_{\mathcal E}''(s)\rVert^2
=\kappa_g^2(s)+\kappa_n^2(s).
\]

Die geodätische Krümmung \(\kappa_g\) wird benötigt, weil sie die
vorzeichenbehaftete Richtungsänderung innerhalb der Referenzfläche misst. Die
Normalkrümmung \(\kappa_n\) wird benötigt, weil selbst ein innerhalb der
Fläche gerader Pfad sich im dreidimensionalen Raum krümmt, während sich die
Erde darunter krümmt. Die gesamte räumliche Krümmung allein würde diese beiden
ingenieurtechnischen Bedeutungen vermischen.

\section{Welche Krümmung trägt die Alignment-Änderung?}

Für die Flächenrealisierung ist die eisenbahnrelevante horizontale
Drehkrümmung
\[
\kappa_{\mathrm{rail}}(s)=\kappa_g(s).
\]
Das bearbeitete intrinsische Krümmungsgesetz wird somit vorbehaltlich des
gewählten Flächenkontexts und der Randdaten als erforderlicher Verlauf der
geodätischen Krümmung realisiert. Normalkrümmung ist eine Folge dieses
Kontexts; sie ist keine zusätzliche horizontale Alignment-Änderung.

In der ebenen Realisierung ist die Flächennormale konstant und
\(\kappa_n=0\). Die Beziehung reduziert sich lokal auf
\[
\kappa_g(s)=\kappa(s)=\theta'(s).
\]
Die ebene Ingenieurrechnung wird daher nicht ungültig. Sie ist die anwendbare
lokale Realisierung, wenn Flächen- und Projektionseffekte unter der von der
Aufgabe geforderten Genauigkeit liegen. Ellipsoidische Realisierung wird
wichtig, wenn Ausdehnung, geforderte Genauigkeit, Projektionsverzerrung oder
erdbezogene räumliche Folgen diese Effekte wesentlich machen.

\section{Projektion ist eine Darstellung des Ergebnisses}

Eine projizierte Kurve \(C(\gamma_{\mathcal E})\) unterstützt CAD, GIS,
Austausch und Berechnung. Direkt aus dieser Repräsentation berechnete
Krümmung kann von \(\kappa_g\) abweichen, weil die Projektion metrische
Eigenschaften verändert. Das macht projizierte Arbeit nicht unbrauchbar; es
macht ihre Anwendbarkeit und Genauigkeit zu Bestandteilen des
Ausführungsvertrags.

Die Erdrealisierung empfängt den intrinsischen Zustand, die Referenzfläche,
Metrik, Orientierung und Randdaten. Sie erzeugt eine Flächentrajektorie und
deren geodätische/normale Zerlegung. Sie bewahrt Alignment-Identität und
Änderungsprovenienz. Sie kann weder ein Projekt-CRS auswählen, noch beweisen,
dass jede Aufgabe eine ellipsoidische Behandlung erfordert, oder entscheiden,
ob das Ergebnis akzeptabel ist.

Das nächste Kapitel verwendet genau diese Ausgaben, um die räumliche Pose zu
konstruieren und Krümmung, Überhöhung, Gravitation und Geschwindigkeit zu
verbinden.
